\documentclass[../ap_2013.tex]{subfiles}

\graphicspath{{./image/}}

\begin{document}

\setcounter{chapter}{2}
\chapter{複素解析}
\section{}
\subsection{}
\begin{equation}\begin{aligned}[b]
    f(z) &= \frac{1}{z-1}\qty(\frac{1}{2}\frac{1}{z}-\frac{1}{2}\frac{1}{z+2})\\
     &= \frac{1}{z-1}\qty(\frac{1}{2}\frac{\frac{1}{z-1}}{1+\frac{1}{z-1}}-\frac{1}{6}\frac{1}{1+(z-1)/3})\\
     &= \frac{1}{z-1}\qty[\frac{1}{2}\frac{1}{z-1}\sum_{n=0}^{\infty}\frac{(-1)^n}{(z-1)^n}-\frac{1}{6}\sum_{n=0}^{\infty}\qty(-\frac{1}{3})^n(z-1)^n]\\
     &= \frac{1}{2}\sum_{n=-\infty}^{-2}\frac{1}{(z-1)^n}+\frac{1}{2}\sum_{n=-1}^{\infty}\qty(-\frac{1}{3})^n(z-1)^n
\end{aligned}\end{equation}

\subsection{}
留数定理より
\begin{equation}\begin{aligned}[b]
    \int_C f(z)dz &= 2\pi i \Res[f,z=1] = -\frac{\pi i}{3}
\end{aligned}\end{equation}

\section{}
\subsection{}
\begin{equation}\begin{aligned}[b]
    z^4+1&=0\\
    z &= e^{i\pi/4},\,e^{i3\pi/4},\,e^{i5\pi/4},\,e^{i7\pi/4}
\end{aligned}\end{equation}
なので、上半平面にある極は
\begin{equation}\begin{aligned}[b]
    z =e^{i\pi/4},\,e^{i3\pi/4}
\end{aligned}\end{equation}

\subsection{}
積分経路\(C\)を\((-R,R)\)の実軸上\(C_1(R)\)と半径\(R\)の上半円の弧\(C_2(R)\)を合わせたもの、\(R\to \infty\)としたものを考える。
\begin{equation}\begin{aligned}[b]
    \int_C g(z) dz &= \lim_{R\to\infty}\int_{C_1(R)}g(z)dz + \int_{C_2(R)}g(z)dz
\end{aligned}\end{equation}
これの左辺は留数定理より
\begin{equation}\begin{aligned}[b]
    \int_C g(z)dz &= 2\pi i\Res[g,z=e^{i\pi/4}]+2\pi i\Res[g,z=e^{i3\pi/4}]\\
    &=\frac{\pi}{2}e^{i\pi/4}+\frac{\pi}{2}e^{-i\pi/4}
    =\frac{\pi}{\sqrt{2}}
\end{aligned}\end{equation}
右辺の円弧の経路は\(z=Re^{i\theta}\)として、
\begin{equation}\begin{aligned}[b]
    \int_{C_2(R)}g(z)dz
    &= \int_{0}^{\pi}g(Re^{\theta})iRe^{i\theta}\\
    &= \int_{0}^{\pi}\frac{R^2e^{2i\theta}}{R^4e^{4i\theta}+1}iRe^{i\theta}d\theta\\
    &\leq \int_{0}^{\pi} \frac{d\theta}{R} \to 0 \qquad (R\to\infty)
\end{aligned}\end{equation}
となる。
よって
\begin{equation}\begin{aligned}[b]
    \int_C g(z) dz &= \lim_{R\to\infty}\int_{C_1(R)}g(z)dz \\
    \int_{-\infty}^{\infty} \frac{x^2}{x^4+1}dx&= \frac{\pi}{\sqrt{2}}
\end{aligned}\end{equation}



\end{document}
